% !TeX root = ../main.tex
论文写完，也意味着研究生阶段快要结束了。回顾这段时间，从确定选题到查阅资料，从方法设计到实验实现，再到后面的论文撰写和反复修改，整个过程并不轻松。中间遇到过很多具体问题，也走过一些弯路。能够比较顺利地完成这篇论文，离不开很多老师、同学和家人的帮助与支持。

首先，我要感谢我的导师王飞副教授。在论文研究和写作过程中，导师在选题确定、思路梳理、实验设计以及论文修改等方面都给予了我很多指导。每当我在方法理解、实验结果分析或者论文表达上遇到问题时，导师都会耐心指出关键问题，并给出明确的修改建议。导师严谨认真的治学态度和踏实负责的工作方式，也让我在这个过程中学到了很多，这些不仅体现在论文完成上，也会对我之后的学习和工作产生影响。

同时，感谢课题组的老师和同学们在整个研究过程中的帮助。平时在实验、代码、论文写作等方面，大家都给过我很多有价值的建议。很多问题其实都是在不断交流和讨论中逐步理清的。和大家一起学习、讨论和推进工作，也让我在研究生阶段收获了很多。

感谢学校和学院提供的学习与科研环境，也感谢研究生阶段各位任课老师的培养。课程学习和日常训练为我后续开展论文工作打下了基础，也让我对本专业的相关知识有了更系统的理解。

最后，感谢我的家人。感谢你们一直以来对我的理解和支持。在论文写作和实验推进比较紧张的时候，正是因为有你们的鼓励和陪伴，我才能够更加安心地完成研究生阶段的学习任务，并坚持把这篇论文完成下来。

在此，向所有在我研究生阶段给予帮助和支持的老师、同学和家人表示衷心的感谢。
